\chapter{Matrix Product Vectors}\label{ch:mps}

Fix a physical dimension $d$ and a bond dimension $D$.
This chapter studies the tensors that generate matrix product vector
(MPV) families.
All tensors here are translation-invariant with periodic
boundary conditions (PBC).

\section{Basic definitions}

\begin{definition}[MPS tensor]\label{def:mps_tensor}
    \lean{MPSTensor}
    \leanok
    A (translation-invariant, PBC) \emph{MPS tensor} with
    physical dimension~$d$ and bond dimension~$D$ is a collection of matrices
    $\{A^i\}_{i=0}^{d-1}$, where $A^i \in \MN{D}$, indexed by a physical index
    $i \in \{0, \ldots, d{-}1\}$.
    Such a tensor defines an MPV family. Diagrammatically,
    \begin{center}
        % A^i \in \MN{D}: the node is A, its upper leg is the physical
        % index i, and the open west/east legs are the two virtual
        % matrix indices. Boundary signature: (virtual-west=1 |
        % virtual-east=1 | physical-up=1 | physical-down=0).
        \tenkzkernel
        \begin{tenkz}[cols=1, boundary=open]
            \tn[ports={90:physical:$i$}]{A}
        \end{tenkz}
    \end{center}
    The black node denotes the tensor $A$, the horizontal legs are virtual,
    and the upper leg is the physical index $i$.
\end{definition}

\begin{definition}[Word evaluation]\label{def:eval_word}
    \lean{Kraus.evalWord}
    \leanok
    \uses{def:mps_tensor}
    Given a word $w = (i_1, \ldots, i_L) \in \{0,\ldots,d{-}1\}^L$,
    the \emph{word evaluation} is the matrix product
    \begin{align}
        A^w
         & := A^{i_1} A^{i_2} \cdots A^{i_L} \in \MN{D}.
        \label{eq:mps_word_eval}
    \end{align}
    The empty word evaluates to the identity, $A^\varnothing = \Id_D$.
    In tensor-network notation,
    \begin{center}
        % A^w = A^{i_1} A^{i_2} \cdots A^{i_L}: horizontal bonds are
        % contracted matrix indices, the two outer virtual legs remain
        % open, and the upper legs carry the letters of w. The ellipsis
        % and brace make the L-site word schematic. The ink therefore
        % displays three representative physical legs; the ellipsis carries
        % the remaining L-3 factors. Emitted boundary signature of this
        % finite schematic: (virtual-west=1 | virtual-east=1 |
        % physical-up=3 | physical-down=0).
        \tenkzkernel
        \begin{tenkz}[cols=4, boundary=open, physical=up]
            \tn[label pos=135, ports={90:physical:$i_1$}]{A} &
            \tn[label pos=135]{A} & \tn[skin=dots]{} &
            \tn[label pos=135, ports={90:physical:$i_L$}]{A}
            \tnmark[form=bracket]{(1,1) .. (1,4)}{$L\text{ copies}$}
        \end{tenkz}
    \end{center}
    in which each black node denotes the same local tensor $A$, the virtual
    legs remain open, and the physical legs are labelled by the word
    $(i_1,\ldots,i_L)$.
\end{definition}

\begin{lemma}[Word evaluation along a listed configuration]
    \label{lem:eval_word_of_fn_prod}
    \lean{MPSTensor.evalWord_ofFn_eq_prod}
    \leanok
    \uses{def:eval_word}
    For a configuration $\sigma=(\sigma_0,\ldots,\sigma_{N-1})$, let
    $w_\sigma=[\sigma_0,\ldots,\sigma_{N-1}]$. Then
    $A^{w_\sigma}
        =A^{\sigma_0}A^{\sigma_1}\cdots A^{\sigma_{N-1}}$.
\end{lemma}

\begin{proof}
    \leanok
    \uses{def:eval_word}
    This follows by induction on $N$ from the recursive definition of word
    evaluation.
\end{proof}

\begin{lemma}[Multiplicativity of word evaluation]\label{lem:eval_word_append}
    \lean{Kraus.evalWord_append}
    \leanok
    \uses{def:eval_word}
    For any words $w_1, w_2$,
    $A^{w_1 \cdot w_2} = A^{w_1} A^{w_2}$.
\end{lemma}

\begin{proof}\leanok
    \uses{def:eval_word}
    The claim follows from the definition of $A^w$.
\end{proof}

\begin{lemma}[Centre of $\MN{D}$]\label{lem:center_scalar}
    \lean{Matrix.isScalar_of_commute_span_eq_top}
    \leanok
    If a matrix $Z \in \MN{D}$ commutes with every element of $\MN{D}$, then
    there exists $\lambda \in \C$ such that $Z = \lambda \Id_D$.
\end{lemma}

\begin{proof}\leanok
    Since $Z$ commutes with every matrix unit, all off-diagonal entries of $Z$
    vanish and all diagonal entries are equal.
\end{proof}

\begin{lemma}[Rectangular amplified commutant]
    \label{lem:amplified_center_scalar}
    \lean{Matrix.exists_eq_kronecker_one_of_intertwines_span_eq_top}
    \leanok
    \uses{lem:center_scalar}
    Let $(A_i)_{i\in I}$ span $\MN{D}$, and let
    $C:\C^{m_2}\otimes\C^D\longrightarrow\C^{m_1}\otimes\C^D$ satisfy, for
    every $i\in I$,
    \begin{align}
        (\Id_{m_1}\otimes A_i)C
         & =C(\Id_{m_2}\otimes A_i).
        \notag
    \end{align}
    Then there is a matrix $F:\C^{m_2}\to\C^{m_1}$ such that
    $C=F\otimes\Id_D$.
\end{lemma}

\begin{proof}\leanok
    Each $(a,b)$ multiplicity block of $C$ commutes with every $A_i$.
    Since the $A_i$ span the full matrix algebra,
    Lemma~\ref{lem:center_scalar} makes this block a scalar multiple of the
    identity. These scalars are the entries of $F$.
\end{proof}

\begin{definition}[Matrix product vector]\label{def:mpv}
    \lean{MPSTensor.mpv}
    \leanok
    \uses{def:eval_word}
    The \emph{matrix product vector} (MPV) of a tensor~$A$
    at system size~$N$ is the vector
    \begin{align}
        \ket{V^{(N)}(A)}
         & =
        \sum_{i_1, \ldots, i_N = 0}^{d-1}
        \tr\!(A^{i_1} A^{i_2} \cdots A^{i_N})
        \ket{i_1, \ldots, i_N}
        \in (\C^d)^{\otimes N}.
        \label{eq:mps_mpv_state}
    \end{align}
    Equivalently, for a configuration
    $\sigma = (i_1, \ldots, i_N) \in \{0,\ldots,d{-}1\}^N$, we write
    \begin{align}
        V^{(N)}(A)_\sigma
         & := \tr\!(A^{i_1} A^{i_2} \cdots A^{i_N}).
        \label{eq:mps_mpv_component}
    \end{align}
    The coefficient function
    $\sigma \mapsto V^{(N)}(A)_\sigma$ gives the components of the vector
    in (\ref{eq:mps_mpv_state}); the displayed ket is the corresponding vector
    in $(\C^d)^{\otimes N}$.
    For a general word $w$, we also write $c_w(A) := \tr(A^w)$.
    The \emph{MPV family} generated by $A$ is the collection
    $\mc{V}(A) = \bigl\{\ket{V^{(N)}(A)}\bigr\}_{N \ge 1}$. The coefficient
    $V^{(N)}(A)_\sigma$ is the periodic contraction
    \begin{center}
        % V^{(N)}(A)_{i_1,\ldots,i_N}
        % = \tr(A^{i_1} \cdots A^{i_N}):
        % the periodic return closes the two outer virtual legs into the
        % trace, while the upper legs carry the physical configuration.
        % The ellipsis and brace make the N-site ring schematic. The ink
        % displays three representative physical legs; the ellipsis carries
        % the remaining N-3 factors. Emitted boundary signature of this
        % finite schematic: (virtual-west=0 | virtual-east=0 |
        % physical-up=3 | physical-down=0).
        \tenkzkernel
        \begin{tenkz}[cols=4, boundary=periodic, physical=up]
            \tn[label pos=135, ports={90:physical:$i_1$}]{A} &
            \tn[label pos=135]{A} & \tn[skin=dots]{} &
            \tn[label pos=135, ports={90:physical:$i_N$}]{A}
            \tnmark[form=bracket]{(1,1) .. (1,4)}{$N\text{ copies}$}
        \end{tenkz}
    \end{center}
    of $N$ copies of the local tensor $A$, with the outer virtual legs closed
    by the trace.
\end{definition}

\begin{definition}[Transfer map]\label{def:transfer_map}
    \lean{Kraus.transferMap}
    \leanok
    \uses{def:mps_tensor}
    The \emph{transfer map} associated to a tensor~$A$ is the
    linear map $\E_A : \MN{D} \to \MN{D}$ defined by
    \begin{align}
        \E_A(X)
         & = \sum_{i=0}^{d-1} A^i X (A^i)^\dagger.
        \label{eq:mps_transfer_map}
    \end{align}
    Diagrammatically, the transfer map is the double-layer contraction
    \begin{center}
        % \E_A(X) = \sum_i A^i X (A^i)^\dagger: the upper and conjugate
        % lower nodes are A and its conjugate, their paired physical
        % legs sum i, and X rides the east cup (the input). Traversing
        % the lower wire in the opposite direction supplies the transpose
        % in the adjoint. The west pair remains open as the matrix output.
        % Boundary signature: (virtual-west=2 | virtual-east=0 |
        % physical-up=0 | physical-down=0).
        \tenkzkernel
        \begin{tenkz}[rows={ket,bra}, cols=1, west=open, east={cup=$X$}]
            \tn[label pos=n]{A} \\
            \tn[label pos=s, conjugate]{\overline{A}}
        \end{tenkz}
    \end{center}
    in which the upper node denotes $A$, the lower node denotes $A^\dagger$,
    and the physical index is summed over between them.
\end{definition}

\begin{remark}\leanok
    The transfer map is automatically positive: if $X \ge 0$ then
    $\E_A(X) \ge 0$, because each summand $A^i X (A^i)^\dagger$ is
    positive semidefinite. The abstract framework is developed in the companion quantum-channel
    volume~\cite[``Positive map'']{QICLean2026Blueprint}.
\end{remark}

\section{Site-periodic tensor families}

\begin{definition}[Site-periodic tensor family]\label{def:periodic_mps_tensor}
    \lean{MPSChainTensor}
    \leanok
    For a fixed period $m$, a \emph{site-periodic tensor family} is a
    family of local tensors indexed by $\{0,\ldots,m{-}1\}$, with
    local physical dimension~$d$ and bond dimension~$D$.
\end{definition}

\begin{definition}[Periodic same-state and gauge-equivalence relations]\label{def:periodic_relations}
    \lean{MPSChainTensor.SameState}
    \lean{MPSChainTensor.GaugeEquiv}
    \leanok
    \uses{def:periodic_mps_tensor}
    For site-periodic tensors $A,B$ at fixed period~$m$:
    \begin{itemize}
        \item equality of the induced site-periodic states;
        \item cyclic gauge equivalence.
    \end{itemize}
\end{definition}

\begin{lemma}[Equivalence structures for periodic relations]\label{lem:periodic_relation_equivalences}
    \lean{MPSChainTensor.instEquivalenceSameState}
    \lean{MPSChainTensor.instEquivalenceGaugeEquiv}
    \leanok
    \uses{def:periodic_relations}
    Both periodic relations are equivalence relations (reflexive,
    symmetric, transitive).
\end{lemma}

\section{Gauge equivalence and same MPV}

\begin{definition}[Gauge equivalence]\label{def:gauge_equiv}
    \lean{MPSTensor.GaugeEquiv}
    \leanok
    \uses{def:mps_tensor}
    Two tensors $A$ and $B$ of the same bond dimension~$D$ are
    \emph{gauge equivalent} if there exists an invertible matrix
    $X \in \GL_D(\C)$ such that, for every $i \in \{0,\ldots,d{-}1\}$,
    \begin{align}
        B^i
         & = X A^i X^{-1}.
        \label{eq:mps_gauge_equiv}
    \end{align}
    In tensor-network notation, (\ref{eq:mps_gauge_equiv}) is represented by
    \begin{center}
        % Formula/source: Papers/1606.00608/MPDO-22-12-17-2.tex,
        % lines 187--195, draws B^i = X A^i X^{-1} as I_B=XCX.png.
        % Ink-to-index map: i is the free upper physical index; the west
        % and east wires carry the two virtual matrix indices.
        % Contraction set: the two internal horizontal bonds are summed,
        % implementing the products X A^i and A^i X^{-1}.
        % Type verdict: both sides are rank-three local MPS tensors.
        % Boundary signature: virtual-west=1 | virtual-east=1 |
        % physical-up=1 | physical-down=0.
        \tenkzkernel
        \begin{tenkz}[boundary=open]
            \tn[skin=ring]{X} & \tn[ports={90:physical:$i$}]{A} &
            \tn[skin=ring]{X^{-1}}
        \end{tenkz}
    \end{center}
    where the black node denotes the tensor $A$, the upper leg is the physical
    index~$i$, and the two red side nodes denote the gauge matrices acting on
    the virtual legs.
\end{definition}

\begin{definition}[Same MPV]\label{def:same_mpv}
    \lean{MPSTensor.SameMPV}
    \leanok
    \uses{def:mpv}
    Two tensors $A$ and $B$ (of the same bond dimension)
    \emph{generate the same MPV family}
    if $\ket{V^{(N)}(A)} = \ket{V^{(N)}(B)}$ for every system size~$N$.
\end{definition}

\begin{definition}[Same MPV - different bond dimensions]\label{def:same_mpv2}
    \lean{MPSTensor.SameMPV₂}
    \leanok
    \uses{def:mpv}
    For tensors $A$ and $B$ of possibly different bond
    dimensions $D_1$ and $D_2$,
    we write $\mc{V}(A) = \mc{V}(B)$ if
    $\ket{V^{(N)}(A)} = \ket{V^{(N)}(B)}$ for every word length $N$,
    including the empty word. This algebraic relation is used for exact
    direct-sum identities.
\end{definition}

\begin{definition}[Positive-length same MPV - different bond dimensions]%
    \label{def:same_mpv2_pos}
    \lean{MPSTensor.SameMPV₂Pos}
    \leanok
    \uses{def:mpv}
    For tensors $A$ and $B$ of possibly different bond dimensions,
    we write $\mc{V}^+(A) = \mc{V}^+(B)$ if
    $\ket{V^{(N)}(A)} = \ket{V^{(N)}(B)}$ for every $N > 0$.
\end{definition}

\begin{definition}[Rectangular MPS reduction]
    \label{def:mps_rectangular_reduction}
    \lean{MPSTensor.IsReduction}
    \leanok
    \uses{def:mps_tensor, def:eval_word}
    Let $A$ and $B$ have the same physical alphabet and respective bond
    dimensions $D_A$ and $D_B$.  A reduction from $B$ to $A$ is a pair
    $V\in\C^{D_A\times D_B}$ and
    $W\in\C^{D_B\times D_A}$ such that
    \begin{align}
        VW              & =1_{D_A},
        \label{eq:mps_rectangular_reduction_retraction} \\
        VB^{\mathbf a}W & =A^{\mathbf a}.
        \label{eq:mps_rectangular_reduction_word}
    \end{align}
    for every word $\mathbf a$, including the empty word.
    This is the reduction of
    \cite[Proposition~20 and the following definition]{Molnar2018SemiInjective}.
\end{definition}

\begin{theorem}[Basic rectangular reduction identities]
    \label{thm:mps_rectangular_reduction_identities}
    \lean{MPSTensor.IsReduction.mul_eq_one,
        MPSTensor.IsReduction.evalWord,
        MPSTensor.IsReduction.evalWord_nil,
        MPSTensor.IsReduction.iff_forall_evalWord}
    \leanok
    \uses{def:mps_rectangular_reduction}
    A rectangular reduction satisfies the retraction identity
    (\ref{eq:mps_rectangular_reduction_retraction}) and the word identity
    (\ref{eq:mps_rectangular_reduction_word}).  Conversely, the all-word
    identity alone implies $VW=1_{D_A}$ by taking the empty word.
\end{theorem}

\begin{proof}
    \leanok
    The forward statements are the two defining clauses.  For the converse,
    $A^{\varnothing}=1_{D_A}$ and $B^{\varnothing}=1_{D_B}$, so the empty-word
    instance reads $V1_{D_B}W=1_{D_A}$.
\end{proof}

\begin{remark}\leanok
    Definition~\ref{def:same_mpv} is the same-bond-dimension
    specialization of Definition~\ref{def:same_mpv2}.
    We keep both because later results use both the
    same-dimension and different-dimension forms.
\end{remark}

\begin{definition}[Nonzero proportional MPV]\label{def:nonzero_proportional_mpv}
    \lean{MPSTensor.NonzeroProportionalMPV₂}
    \leanok
    \uses{def:mpv}
    Two tensors $A$ and $B$ (of possibly different bond dimensions)
    generate nonzero proportional MPV families if for every~$N>0$ there is a
    nonzero scalar $c_N \in \C$ such that
    $\ket{V^{(N)}(A)} = c_N \ket{V^{(N)}(B)}$.
    This is the projective reading of the proportionality hypothesis in
    \cite[Theorem~2.10]{Cirac2016MPDO_arXiv}.
    % Maintainer note: projective reading documented in
    % docs/paper-gaps/cpsv16_nonzero_proportionality_reading.tex.
\end{definition}

\begin{definition}[Eventually nonzero proportional MPV]%
    \label{def:eventually_nonzero_proportional_mpv}
    \lean{MPSTensor.EventuallyNonzeroProportionalMPV₂}
    \leanok
    \uses{def:mpv}
    Two tensors $A$ and $B$ are eventually nonzero proportional if, for all
    sufficiently large $N$, there is a nonzero scalar $c_N \in \C$ such that
    $\ket{V^{(N)}(A)} = c_N \ket{V^{(N)}(B)}$. This auxiliary form is used with
    the eventual linear-independence statement
    Lemma~\ref{lem:bnt_finite_index_overlap_orthonormal_li}.
\end{definition}

\begin{lemma}[Elementary nonzero proportionality rules]%
    \label{lem:nonzero_proportional_mpv_basic}
    \lean{MPSTensor.NonzeroProportionalMPV₂.eventually,
        MPSTensor.NonzeroProportionalMPV₂.symm,
        MPSTensor.SameMPV₂Pos.toNonzeroProportionalMPV₂}
    \leanok
    \uses{def:nonzero_proportional_mpv, def:eventually_nonzero_proportional_mpv, def:same_mpv2_pos}
    Nonzero proportionality implies eventual nonzero
    proportionality, is symmetric, and equality of MPV families is the special
    case with scalar $1$, at every positive length.
\end{lemma}

\begin{proof}
    \leanok
    The eventual form follows by viewing an all-length statement as an
    eventual one.
    Symmetry follows by inverting the scalar $c_N$ at each length, and equality
    is the case $c_N = 1$.
\end{proof}

\begin{lemma}[Coefficient isolation after complementary matching]
    \label{lem:finite_coefficient_isolation}
    \lean{LinearIndependent.coefficient_eq_zero_of_sum_eq_of_complement_smul}
    \leanok
    Let $R$ be a ring, let $V$ be an $R$-module, and let $I$ and $K$ be
    finite index sets. Let $T\subseteq I$, and let
    $\{u_i\}_{i\in I}$ and $\{v_k\}_{k\in K}$ be vectors in $V$. Suppose
    that, for every $i\notin T$, there are a scalar $\alpha_i$ and an index
    $\phi(i)\in K$ such that $u_i=\alpha_i v_{\phi(i)}$. If
    \begin{align}
        \sum_{i\in I}a_i u_i=\sum_{k\in K}b_k v_k
        \notag
    \end{align}
    and the coproduct-indexed family $w\colon T\amalg K\to V$, defined by
    $w(\operatorname{inl}(i))=u_i$ and
    $w(\operatorname{inr}(k))=v_k$, is linearly independent, then
    $a_i=0$ for every $i\in T$.
\end{lemma}
\begin{proof}\leanok
    Substitute $u_i=\alpha_i v_{\phi(i)}$ for $i\notin T$ and collect the
    resulting coefficients over the fibres of $\phi$. The given equality is
    then a linear relation in the combined family. Linear independence makes
    every coefficient vanish, in particular each $a_i$ with $i\in T$.
\end{proof}

\begin{lemma}[Scaling of word evaluation]\label{lem:eval_word_smul}
    \lean{Kraus.evalWord_smul}
    \leanok
    \uses{def:eval_word}
    Let $\zeta A$ denote the tensor with matrices
    $\{\zeta A^i\}_{i=0}^{d-1}$. Then, for any word $w$,
    $(\zeta A)^w = \zeta^{|w|} A^w$, where $|w|$ is the length of~$w$.
\end{lemma}

\begin{proof}\leanok
    \uses{def:eval_word}
    Induction on $w$: the base case gives $\zeta^0 = 1$,
    and each step picks up one factor of $\zeta$.
\end{proof}

\begin{lemma}[Scaling of MPV components]\label{lem:mpv_smul}
    \lean{MPSTensor.mpv_smul}
    \leanok
    \uses{def:mpv, lem:eval_word_smul}
    Let $\zeta A$ denote the tensor with matrices
    $\{\zeta A^i\}_{i=0}^{d-1}$. Then, for every $N$ and every configuration
    $\sigma = (i_1, \ldots, i_N)$,
    $V^{(N)}(\zeta A)_\sigma = \zeta^{N} V^{(N)}(A)_\sigma$.
\end{lemma}

\begin{proof}\leanok
    \uses{lem:eval_word_smul}
    Apply Lemma~\ref{lem:eval_word_smul} to the word
    $w = (i_1, \ldots, i_N)$ of length $N$, then take the trace; the scalar
    $\zeta^{N}$ pulls out by linearity.
\end{proof}

\begin{definition}[Gauge-phase equivalence]\label{def:gauge_phase_equiv}
    \lean{MPSTensor.GaugePhaseEquiv}
    \lean{MPSTensor.GaugeEquiv.toGaugePhaseEquiv}
    \leanok
    \uses{def:mps_tensor}
    Two tensors $A$ and $B$ are \emph{gauge-phase equivalent}
    if there exist $X \in \GL_D(\C)$ and $\zeta \in \C \setminus \{0\}$
    such that, for every physical index~$i$,
    \begin{align}
        B^i
         & = \zeta X A^i X^{-1}.
        \label{eq:mps_gauge_phase_equiv}
    \end{align}
\end{definition}

\begin{proof}\leanok
    Gauge equivalence is the special case $\zeta = 1$.
\end{proof}

\begin{definition}[Unit gauge-phase equivalence]
    \label{def:unit_gauge_phase_equiv}
    \lean{MPSTensor.UnitGaugePhaseEquiv}
    \leanok
    \uses{def:gauge_phase_equiv}
    Two tensors are unit gauge-phase equivalent when the gauge-phase equation
    holds with $|\zeta|=1$.
\end{definition}

\begin{lemma}[Unit gauge phase implies gauge phase]
    \label{lem:unit_gauge_phase_to_gauge_phase}
    \lean{MPSTensor.UnitGaugePhaseEquiv.toGaugePhaseEquiv}
    \leanok
    \uses{def:unit_gauge_phase_equiv, def:gauge_phase_equiv}
    Unit gauge-phase equivalence implies gauge-phase equivalence.
\end{lemma}

\begin{proof}\leanok
    A unit-modulus scalar is nonzero.
\end{proof}

\begin{lemma}[Gauge-phase equivalence across equal bond dimensions]
    \label{lem:gauge_phase_equiv_cast_compose}
    \lean{MPSTensor.gaugePhaseEquiv_swap_cast,
        MPSTensor.gaugePhaseEquiv_cast_compose_via_centre}
    \leanok
    \uses{def:gauge_phase_equiv}
    Gauge-phase equivalence is symmetric after identifying equal bond
    dimensions. Moreover, suppose that $A$, $B$, and $C$ have identified bond
    dimensions, that $X_1$ and $X_2$ are invertible, that
    $\zeta_1,\zeta_2\ne0$, and that, for every physical index $i$,
    \begin{align}
        B^i
         & =\zeta_1X_1A^iX_1^{-1},\notag \\
        C^i
         & =\zeta_2X_2A^iX_2^{-1}.
        \notag
    \end{align}
    Then
    \begin{align}
        C^i
         & =(\zeta_2\zeta_1^{-1})(X_2X_1^{-1})
        B^i(X_1X_2^{-1}),
        \label{eq:mps_gauge_phase_compose}
    \end{align}
    so $B$ and $C$ are gauge-phase equivalent after the corresponding
    identification of their bond dimensions.
\end{lemma}

\begin{proof}\leanok
    Symmetry follows by replacing $(X,\zeta)$ with
    $(X^{-1},\zeta^{-1})$. Formula (\ref{eq:mps_gauge_phase_compose}) follows
    by substituting the first gauge formula into the second.
\end{proof}

\begin{lemma}[Gauge covariance of word evaluation]\label{lem:eval_word_gauge}
    \lean{MPSTensor.evalWord_gauge}
    \leanok
    \uses{def:eval_word}
    If $B^i = X A^i X^{-1}$ for all~$i$, then
    $B^w = X A^w X^{-1}$ for every word~$w$.
\end{lemma}

\begin{proof}\leanok
    \uses{def:eval_word}
    Induction on the word length, using $X X^{-1} = \Id$.
\end{proof}

\begin{theorem}[Gauge invariance of MPVs]\label{thm:gauge_invariance}
    \lean{MPSTensor.GaugeEquiv.sameMPV}
    \leanok
    \uses{def:gauge_equiv, def:same_mpv}
    If $A$ and $B$ are gauge equivalent, then they generate the same
    MPV family.
\end{theorem}

\begin{proof}\leanok
    \uses{lem:eval_word_gauge}
    Let $B^i = X A^i X^{-1}$.
    By Lemma~\ref{lem:eval_word_gauge}, $B^w = X A^w X^{-1}$
    for every word~$w$.
    Then $\tr(B^w) = \tr(X A^w X^{-1}) = \tr(A^w)$
    by cyclicity of the trace.
\end{proof}

\begin{theorem}[Gauge-phase scaling of MPVs]\label{thm:gauge_phase_mpv_scaling}
    \lean{MPSTensor.mpv_eq_pow_mul_of_gaugePhase}
    \leanok
    \uses{def:gauge_phase_equiv, def:mpv}
    If $B^i = \zeta X A^i X^{-1}$ for every physical index~$i$, then for
    every system size~$N$ and configuration~$\sigma$ one has
    $V^{(N)}(B)_\sigma = \zeta^N V^{(N)}(A)_\sigma$.
\end{theorem}

\begin{proof}\leanok\uses{lem:eval_word_smul, lem:eval_word_gauge}
    Combine gauge covariance of word evaluation with scalar scaling,
    then take traces.
\end{proof}

\section{Injectivity and normality}

\begin{definition}[Injectivity]\label{def:injective}
    \lean{Kraus.IsInjective}
    \leanok
    \uses{def:mps_tensor}
    A tensor~$A$ is \emph{injective} if its matrices span the
    full matrix algebra:
    $\spn_{\C}\{A^i : i = 0,\ldots,d{-}1\} = \MN{D}$.
\end{definition}

\begin{definition}[Coefficient-dual inverse and reduction cross matrix]
    \label{def:mps_reduction_cross_matrix}
    \lean{MPSTensor.coefficientDualInverse,
        MPSTensor.reductionCrossMatrix}
    \leanok
    \uses{def:mps_tensor, def:injective}
    Let $\widetilde A=(\widetilde A^i)_i$ be injective, with bond dimension
    $D_A$, and choose the coefficient dual $C=(C^i)_i$ from a right inverse
    to the linear-combination map of $\widetilde A$.  Thus
    $C^i_{a'a}$ is the coefficient of $\widetilde A^i$ in the chosen
    expansion of the matrix unit $\lvert a\rangle\!\langle a'\rvert$.
    For a tensor $\widetilde B$ of bond dimension $D_B$, define the
    mixed-bond matrix on $\C^{D_B}\otimes\C^{D_A}$ by
    \begin{align}
        K(\widetilde B,C)
         & =\sum_i \widetilde B^i\otimes(C^i)^T.
        \notag
    \end{align}
    The superscript $T$ is the ordinary transpose, not the conjugate
    transpose.  This is the matrix drawn in
    \cite[proof of Proposition~20]{Molnar2018SemiInjective},
    \texttt{cornerproblem.tex} lines 3817--3838.
\end{definition}

\begin{lemma}[Coefficient dual and reversed-word power formulas]
    \label{lem:mps_reduction_cross_matrix_power}
    \lean{MPSTensor.coefficientDualInverse_sum_entry,
        MPSTensor.reductionCrossMatrix_apply,
        MPSTensor.reductionCrossMatrix_pow_eq_sum,
        MPSTensor.reductionCrossMatrix_pow_apply,
        MPSTensor.reductionCrossMatrix_coefficientDualInverse}
    \leanok
    \uses{def:mps_reduction_cross_matrix}
    The coefficient dual obeys the exact matrix-unit identity
    \begin{align}
        \sum_i \widetilde A^i_{xy}C^i_{a'a}
         & =\delta_{x,a}\delta_{y,a'}.
        \notag
    \end{align}
    Consequently,
    \begin{align}
        (K(\widetilde B,C)^n)_{(b,a),(b',a')}
         & =\sum_{i_1,\ldots,i_n}
        (\widetilde B^{i_1}\cdots\widetilde B^{i_n})_{bb'}
        (C^{i_n}\cdots C^{i_1})_{a'a}.
        \notag
    \end{align}
    In particular, pairing $\widetilde A$ with its own dual gives the
    rank-one matrix whose entry is
    $\delta_{x,a}\delta_{y,a'}$.  The reversed order
    $C^{i_n}\cdots C^{i_1}$ is forced by the transpose and records the
    upper-leg orientation in
    \texttt{cornerproblem.tex} lines 3817--3865.
\end{lemma}

\begin{proof}
    \leanok
    \uses{def:mps_reduction_cross_matrix}
    The coefficient identity follows by applying the chosen right inverse to a
    matrix unit.  Expanding the power of the finite sum defining $K$,
    multiplicativity of the Kronecker product collects the lower factors in
    site order.  Transposing the upper product reverses its factors, which gives
    the displayed entry formula.  Specializing the lower tensor to
    $\widetilde A$ and applying the coefficient identity gives the stated
    rank-one matrix.
\end{proof}

\begin{theorem}[Cross trace powers and source open-boundary contraction]
    \label{thm:mps_reduction_cross_trace_open}
    \lean{MPSTensor.SameMPV₂Pos.trace_evalWord,
        MPSTensor.trace_reductionCrossMatrix_pow_of_sameMPV₂Pos,
        MPSTensor.reductionOpenBoundaryMatrixContraction,
        MPSTensor.reductionOpenBoundaryContraction,
        MPSTensor.reductionOpenBoundaryMatrixContraction_eq_cross_sum,
        MPSTensor.reductionOpenBoundaryMatrixContraction_self,
        MPSTensor.reductionOpenBoundaryContraction_self,
        MPSTensor.reductionOpenBoundaryContraction_of_sameMPV₂Pos,
        MPSTensor.reductionOpenBoundaryMatrixContraction_blockTensor_of_sameMPV₂Pos}
    \leanok
    \uses{def:same_mpv2_pos, def:injective,
        def:mps_reduction_cross_matrix, thm:same_mpv2_pos_block_tensor}
    Suppose $\widetilde A$ and $\widetilde B$ have equal positive-length
    MPV coefficients and $\widetilde A$ is injective.  This equality hypothesis
    is the $\lambda=1$ specialization of the proportional premise printed in
    Proposition~20; the unrestricted unscaled conclusion is false
    \cite{gap:mgsc18_reduction_proportionality_scalar}.  Equality transports
    the trace of every nonempty word.  For every $n>0$,
    \begin{align}
        \tr K(\widetilde B,C)^n
         & =D_A^n.
        \notag
    \end{align}

    For arbitrary tensors $B_0,C_0$, every $n\geq0$, and every boundary matrix
    $X$, define
    \begin{align}
        \mathcal O_n(B_0,C_0;X)_{a'a}
         & =\sum_{i_1,\ldots,i_n}
        \tr\!\left(B_0^{i_1}\cdots B_0^{i_n}X\right)
        (C_0^{i_n}\cdots C_0^{i_1})_{a'a}.
        \notag
    \end{align}
    If $B_0$ has bond dimension $D_B$ and $C_0$ has bond dimension $D_A$,
    its exact cross-sum expansion is
    \begin{align}
        \mathcal O_n(B_0,C_0;X)_{a'a}
         & =\sum_{x,y}K(B_0,C_0)^n_{(x,a),(y,a')}X_{yx}.
        \notag
    \end{align}
    For $n>0$, if $B_0=\widetilde A$, $C_0=C$, and
    $X\in\C^{D_A\times D_A}$, then
    \begin{align}
        \mathcal O_n(\widetilde A,C;X)_{a'a}
         & =D_A^{n-1}X_{a'a}.
        \notag
    \end{align}
    The same-alphabet open-boundary contraction is the specialization
    $X=B_0^w$.  Thus $X=\widetilde A^w$ gives the target self-evaluation, while
    equality of the positive-length MPV coefficients gives
    \begin{align}
        \mathcal O_n(\widetilde B,C;\widetilde B^w)_{a'a}
         & =D_A^{n-1}(\widetilde A^w)_{a'a}.
        \notag
    \end{align}

    For the padded blocked contraction, suppose additionally that the original
    tensors $A$ and $B$ have equal positive-length MPV coefficients.  If
    $\widetilde A$ and $\widetilde B$ are obtained by blocking $p>0$ original
    sites, then for every unblocked lower tail word $w$,
    \begin{align}
        \mathcal O_n(\widetilde B,C;B^w)_{a'a}
         & =D_A^{n-1}(A^w)_{a'a}.
        \notag
    \end{align}
    The trace-power identity is the closed contraction at
    \texttt{cornerproblem.tex} lines 3839--3865, and the open-boundary identity
    is the padded contraction at lines 3866--3887.  The source next uses a
    rank-one rewrite and comparison to derive
    $VB^{i_1}\cdots B^{i_m}W=A^{i_1}\cdots A^{i_m}$ at lines 3888--3938.
\end{theorem}

\begin{proof}
    \leanok
    \uses{lem:mps_reduction_cross_matrix_power, lem:eval_word_block}
    Expand $K^n$ by Lemma~\ref{lem:mps_reduction_cross_matrix_power} and use
    the trace formula for a Kronecker product.  Positive-length MPV equality
    replaces the trace of each nonempty $\widetilde B$ word by the
    corresponding $\widetilde A$ trace.  The resulting scalar sum is therefore
    the trace of the target cross matrix power.  The coefficient-dual identity
    gives that target cross matrix its rank-one form, whose $n$th trace is
    $D_A^n$.

    Expanding the lower trace into matrix entries gives the cross-sum formula
    for arbitrary $X$.  For the target tensor, the rank-one form of
    $K(\widetilde A,C)$ evaluates this sum as $D_A^{n-1}X_{a'a}$.
    Specializing $X$ to a word evaluation gives the self-evaluation formula.

    For blocked tensors, Lemma~\ref{lem:eval_word_block} identifies each
    blocked word with the evaluation of its flattened word.  Since $n,p>0$,
    this flattened word is nonempty.  Append the unblocked tail $w$, apply
    positive-length MPV equality to the resulting original-site word, and use
    the target self-evaluation with $X=A^w$.
\end{proof}

\begin{theorem}[Rank-one factorizations]
    \label{thm:rank_one_factorizations}
    \lean{LinearMap.exists_smulRight_of_finrank_range_eq_one,
        Matrix.exists_eq_vecMulVec_of_rank_eq_one,
        Matrix.exists_rectangular_factors_of_rank_eq_one}
    \leanok
    Let $f\colon X\to Y$ be a linear map over a division ring $K$ whose range
    has dimension one.  Then there are a nonzero linear functional
    $\varphi\colon X\to K$, a nonzero vector $y\in Y$, and a vector $x\in X$
    such that
    \begin{align}
        \varphi(x) & =1, & f(x) & =y, & f(z) & =\varphi(z)y
        \quad\text{for every }z\in X.
        \notag
    \end{align}
    Consequently, if $I$ and $J$ are finite index sets, then every rank-one
    matrix $M\in K^{I\times J}$ over a field is an outer product of two
    nonzero vectors: there are $w\in K^I$ and $v\in K^J$ such that
    \begin{align}
        M_{ij} & =w_i v_j.
        \notag
    \end{align}
    In particular, if both indices of $M$ are the product
    $D_B\times D_A$, then there are nonzero rectangular matrices
    $W\in K^{D_B\times D_A}$ and $V\in K^{D_A\times D_B}$ such that
    \begin{align}
        M_{(b,a),(b',a')}
         & =W_{ba}V_{a'b'}.
        \notag
    \end{align}
    This is the rank-one factorization used in
    \cite[proof of Proposition~20]{Molnar2018SemiInjective},
    \texttt{cornerproblem.tex} lines 3866--3936.
\end{theorem}

\begin{proof}
    \leanok
    Choose an isomorphism from the scalar division ring to the one-dimensional
    range of $f$.  Composing its inverse with $f$ gives $\varphi$, while the
    image of $1$ gives $y$.  Choose a preimage $x$ of $y$; the construction gives
    $\varphi(x)=1$, $f(x)=y$, and the stated factorization of $f$.

    In standard bases this factorization is an outer product of two nonzero
    coordinate vectors.  For the product-index specialization, reshape the
    left vector as a $D_B\times D_A$ matrix and the right vector as a
    $D_A\times D_B$ matrix, preserving the displayed index order.
\end{proof}

The inverse-tensor characterization of injectivity says that $A$ is
injective exactly when there is a tensor $A^{-1}$ such that
\begin{align}
    \sum_i (A^i)_{\alpha,\beta}(A^{-1,i})_{\alpha',\beta'}
     & =
    \delta_{\alpha,\alpha'}\delta_{\beta,\beta'}.
    \label{eq:mps_injective_inverse}
\end{align}
Thus (\ref{eq:mps_injective_inverse}) identifies the two pairs of virtual
indices when $A^{-1}$ is contracted with $A$ through their common physical
index. In shorthand,
% Correct equation, not a figure: the source states the unindexed injective
% identity in prose. I_Injectivity.png illustrates the following, strictly
% stronger block-injective identity with the additional j,j' index pair.
$A^{-1}A=\Id$. This is the identity used in
\cite[Section~2.3, lines~324--331]{Cirac2016MPDO_arXiv}.

\begin{definition}[$L$-block injectivity]\label{def:l_blk_injective}
    \lean{Kraus.IsNBlkInjective}
    \leanok
    \uses{def:mps_tensor}
    A tensor~$A$ is \emph{$L$-block injective} if
    \begin{align}
        \spn_{\C}\{A^{i_1} \cdots A^{i_L}
        : (i_1,\ldots,i_L) \in \{0,\ldots,d{-}1\}^L\}
         & = \MN{D}.
        \label{eq:mps_block_injective}
    \end{align}
    An injective tensor is $1$-block injective.
\end{definition}

\begin{lemma}[One-site injectivity gives one-block injectivity]
    \label{lem:one_block_injective_of_injective}
    \lean{Kraus.isNBlkInjective_one_of_isInjective}
    \leanok
    \uses{def:injective, def:l_blk_injective}
    If the one-site matrices $\{A^i\}_{i=0}^{d-1}$ span $\MN{D}$, then the
    length-one word products span $\MN{D}$, so $A$ is $1$-block injective.
\end{lemma}

\begin{proof}
    \leanok
    The length-one words are exactly the letters $i \in \{0,\ldots,d{-}1\}$, and
    evaluating a length-one word gives the matrix $A^i$. Thus the length-one
    word span is the same span as in one-site injectivity.
\end{proof}

\begin{definition}[Normal tensor]\label{def:normal}
    \lean{Kraus.IsNormal}
    \leanok
    \uses{def:l_blk_injective}
    A tensor~$A$ is \emph{normal}
    \cite{Cirac2021Matrix}
    if there exists a positive integer $L_0$ such that $A$ is $L_0$-block
    injective.
\end{definition}

\begin{theorem}[Equality-case existence of rectangular reductions]
    \label{thm:mps_equality_reduction_exists}
    \lean{MPSTensor.exists_isReduction_of_isInjective_of_sameMPV₂Pos,
        MPSTensor.exists_isReduction_of_isNormal_of_sameMPV₂Pos}
    \leanok
    \uses{def:mps_rectangular_reduction, def:same_mpv2_pos, def:injective,
        def:normal}
    Let $A$ and $B$ have the same physical alphabet and positive target bond
    dimension $D_A>0$.  Suppose that
    $\ket{V^{(N)}(B)}=\ket{V^{(N)}(A)}$ for every $N>0$.
    If $A$ is injective, then there exist
    $V\in\C^{D_A\times D_B}$ and $W\in\C^{D_B\times D_A}$ forming a
    reduction from $B$ to $A$.  The same conclusion holds when $A$ is normal.

    This is the $\lambda=1$ equality case of
    \cite[Proposition~20]{Molnar2018SemiInjective}; see
    \texttt{cornerproblem.tex} lines 3815--3938.  The unrestricted
    proportional premise does not imply the unscaled conclusion
    \cite{gap:mgsc18_reduction_proportionality_scalar}.
\end{theorem}

\begin{proof}
    \leanok
    \uses{lem:one_block_injective_of_injective,
        lem:isNBlkInjective_iff_blockTensor_isInjective,
        thm:same_mpv2_pos_block_tensor,
        thm:mps_reduction_cross_trace_open,
        thm:mpu_semisimple_scaled_trace_rank_one,
        thm:rank_one_factorizations}
    First suppose that $A$ is injective; equivalently, use blocking length one.
    More generally, choose any positive blocking length $p$ for which the
    blocked target $\widetilde A=A^{[p]}$ is injective, and put
    $\widetilde B=B^{[p]}$.  Positive blocking preserves equality of the MPV
    families.  Let $C$ be the coefficient dual of $\widetilde A$ and set
    \begin{align}
        K
         & =\sum_i \widetilde B^i\otimes(C^i)^{\mathsf T}.
        \notag
    \end{align}
    Regard $K$ as an endomorphism and take its Jordan--Chevalley decomposition
    $K=N+S$, where $N$ is nilpotent, $S$ is semisimple, and both belong to the
    algebra generated by $K$.  Their membership in this commutative generated
    algebra implies $NS=SN$.  The commuting nilpotent part does not change
    traces of powers.  Since $\tr K^n=D_A^n$ for every $n>0$, the semisimple
    endomorphism $D_A^{-1}S$ has trace one at every power above one.  The
    semisimple trace-power theorem gives $\rank S=1$; commutation and
    nilpotence then imply $NS=SN=0$ and
    \begin{align}
        K^J
         & =S^J
        \notag
    \end{align}
    for some $J>0$.

    Factor the matrix of $S$ with the mixed-bond orientation
    \begin{align}
        S_{(b,a),(b',a')}
         & =W_{ba}V_{a'b'}.
        \notag
    \end{align}
    Its first trace is $D_A$, so multiplication of outer products gives
    \begin{align}
        S^J
         & =D_A^{J-1}S.
        \notag
    \end{align}
    Apply the padded open-boundary identity with an arbitrary original-scale
    word $w$.  Rewriting its cross sum with the preceding factorization yields
    \begin{align}
        D_A^{J-1}(VB^wW)_{a'a}
         & =D_A^{J-1}(A^w)_{a'a}.
        \notag
    \end{align}
    Since $D_A>0$, cancel the nonzero scalar and conclude
    $VB^wW=A^w$ for every word $w$.  The empty word gives $VW=1_{D_A}$.
    For normal $A$, choose the positive injective blocking supplied by
    normality and repeat the same padded argument; the lower tail remains the
    unblocked word $w$, so the conclusion is already at the original physical
    scale.
\end{proof}


\begin{definition}[Residual algebra of a selected reduction]
    \label{def:mps_reduction_residual_algebra}
    \lean{MPSTensor.reductionResidual,
        MPSTensor.reductionResidualGeneratorSet,
        MPSTensor.reductionResidualAlgebra,
        MPSTensor.reductionResidual_mem_reductionResidualAlgebra}
    \leanok
    \uses{def:mps_rectangular_reduction}
    Let $(V,W)$ be a reduction from $B$ to $A$.  For each physical letter
    $i$, define
    \begin{align}
        N^i & = B^i-WA^iV.
        \label{eq:mps_reduction_residual_letter}
    \end{align}
    The \emph{residual algebra} $\mathcal R(B,A;V,W)$ is the nonunital
    subalgebra of $\MN{D_B}$ generated by the matrices $N^i$.
    Each residual letter belongs to $\mathcal R(B,A;V,W)$.
    This is the residual algebra of Proposition~21 in
    \cite{Molnar2018SemiInjective}.
\end{definition}

\begin{definition}[Residual nilpotency length]
    \label{def:mps_reduction_residual_nilpotency_length}
    \lean{MPSTensor.IsReductionResidualNilpotencyBound,
        MPSTensor.reductionResidualNilpotencyLength}
    \leanok
    \uses{def:mps_reduction_residual_algebra}
    A natural number $L$ is a \emph{residual-word nilpotency bound} if
    every word $N^{i_1}\cdots N^{i_L}$ of exactly $L$ residual letters is
    zero.  Definition~8 of \cite{Molnar2018SemiInjective} requires all
    residual words of every length $n\geq L$ to vanish.  The formulations
    are equivalent: an exact-length-zero prefix forces every longer word to
    vanish.  The residual nilpotency length $\ell(B,A;V,W)$ is the infimum
    of the set of exact-length bounds.  When a bound exists, this is its
    least value and equals the nilpotency length of Definition~8.
\end{definition}

\begin{theorem}[Exterior-word strengthening of the residual sandwich]
    \label{thm:mps_reduction_residual_exterior_sandwich}
    \lean{MPSTensor.IsReduction.evalWord_mul_reductionResidual_mul_evalWord_eq_zero}
    \leanok
    \uses{thm:mps_rectangular_reduction_identities}
    Under the same reduction hypothesis, for arbitrary exterior words
    $\mathbf p,\mathbf q$ and every nonempty word $\mathbf u$,
    \begin{align}
        VB^{\mathbf p}N^{\mathbf u}B^{\mathbf q}W & =0.
        \label{eq:mps_reduction_residual_exterior_sandwich}
    \end{align}
    This buffered statement is a derived TNLean strengthening; it is not the
    statement of Lemma~\texttt{VNW=0}.
\end{theorem}

\begin{proof}
    \leanok
    Apply $VB^{\mathbf w}W=A^{\mathbf w}$ to $\mathbf p$, $\mathbf q$, and
    $\mathbf p\,i\,\mathbf q$.  The two expressions for
    $A^{\mathbf p}A^iA^{\mathbf q}$ give
    \begin{align}
        VB^{\mathbf p}N^iB^{\mathbf q}W & =0.
        \notag
    \end{align}
    Induction on the nonempty residual word $\mathbf u$ proves
    (\ref{eq:mps_reduction_residual_exterior_sandwich}).
\end{proof}

\begin{theorem}[Positive residual words vanish under the reduction]
    \label{thm:mps_reduction_residual_sandwich}
    \lean{MPSTensor.IsReduction.evalWord_reductionResidual_sandwich_eq_zero}
    \leanok
    \uses{def:mps_rectangular_reduction,
        def:mps_reduction_residual_algebra,
        thm:mps_reduction_residual_exterior_sandwich}
    Let $(V,W)$ be a reduction from $B$ to $A$.  For every nonempty word
    $\mathbf u$,
    \begin{align}
        VN^{\mathbf u}W & =0
        \qquad (|\mathbf u|>0).
        \label{eq:mps_reduction_residual_sandwich}
    \end{align}
    No equality of closed chains is required.
    This is Lemma~\texttt{VNW=0} of
    \cite{Molnar2018SemiInjective}; see
    \texttt{cornerproblem.tex} lines 3946--3970.
\end{theorem}

\begin{proof}
    \leanok
    \uses{thm:mps_reduction_residual_exterior_sandwich}
    Take both exterior words empty in
    (\ref{eq:mps_reduction_residual_exterior_sandwich}).
\end{proof}

\begin{definition}[Strict residual-window sums]
    \label{def:mps_reduction_strict_window_sums}
    \lean{MPSTensor.reductionInitialWindowSum,
        MPSTensor.reductionStrictWindowSum,
        MPSTensor.reductionInitialWindowRangeSum,
        MPSTensor.reductionStrictWindowRangeSum}
    \leanok
    \uses{def:mps_rectangular_reduction,
        def:mps_reduction_residual_algebra, def:eval_word}
    Let $(V,W)$ be a reduction from $B$ to $A$.  For a word $\mathbf w$, set
    $\mathcal I(\varnothing)=\mathcal W(\varnothing)=0$ and define
    recursively
    \begin{align}
        \mathcal I(i\mathbf w)
         & =WA^iV\bigl(N^{\mathbf w}+\mathcal I(\mathbf w)\bigr),
        \label{eq:mps_reduction_initial_window_recursion}         \\
        \mathcal W(i\mathbf w)
         & =\mathcal I(i\mathbf w)+N^i\mathcal W(\mathbf w).
        \label{eq:mps_reduction_strict_window_recursion}
    \end{align}
    For $\mathbf i=(i_1,\ldots,i_n)$, also set
    \begin{align}
        \widehat{\mathcal I}(\mathbf i)
         & =\sum_{m=1}^{n}
        W A^{i_1}\cdots A^{i_m}V
        N^{i_{m+1}}\cdots N^{i_n},                                      \\
        \widehat{\mathcal W}(\mathbf i)
         & =\sum_{r=0}^{n-1}\sum_{m=1}^{n-r}
        N^{i_1}\cdots N^{i_r}
        W A^{i_{r+1}}\cdots A^{i_{r+m}}V
        N^{i_{r+m+1}}\cdots N^{i_n}.
        \label{eq:mps_reduction_strict_window_sum}
    \end{align}
    The substitution $s=r+m$ identifies the second indexing set with
    $0\leq r<s\leq n$.  Thus the central $A$-word is nonempty, and
    $\widehat{\mathcal I}$ is the part with $r=0$.
\end{definition}

\begin{theorem}[Recursive and explicit strict-window sums]
    \label{thm:mps_reduction_strict_window_range}
    \lean{MPSTensor.IsReduction.reductionStrictWindowSum_eq_reductionStrictWindowRangeSum}
    \leanok
    \uses{def:mps_reduction_strict_window_sums,
        thm:mps_rectangular_reduction_identities}
    Under the reduction hypothesis, for every word $\mathbf i$,
    \begin{align}
        \mathcal W(\mathbf i)=\widehat{\mathcal W}(\mathbf i).
    \end{align}
\end{theorem}

\begin{proof}
    \leanok
    Induct on the length of the word.  The summands with $r=0$ constitute
    $\widehat{\mathcal I}$, while deleting the first letter identifies the
    summands with $r>0$ with $N^{i_1}\widehat{\mathcal W}(i_2,\ldots,i_n)$.
    The identity $VW=1$ identifies $\widehat{\mathcal I}$ with
    $\mathcal I$.
\end{proof}

\begin{theorem}[Corrected strict-window residual expansion]
    \label{thm:mps_reduction_residual_expansion}
    \lean{MPSTensor.IsReduction.evalWord_eq_reductionResidual_add_strictWindowRangeSum}
    \leanok
    \uses{def:mps_rectangular_reduction,
        def:mps_reduction_strict_window_sums,
        thm:mps_reduction_strict_window_range}
    Let $(V,W)$ be a reduction from $B$ to $A$.  For every word
    $\mathbf i=(i_1,\ldots,i_n)$,
    \begin{align}
        B^{i_1}\cdots B^{i_n}
         & =N^{i_1}\cdots N^{i_n}
        +\sum_{0\leq r<s\leq n}
        N^{i_1}\cdots N^{i_r}
        W A^{i_{r+1}}\cdots A^{i_s}V
        N^{i_{s+1}}\cdots N^{i_n}.
        \label{eq:mps_reduction_residual_expansion}
    \end{align}
    This statement requires only the reduction.  The weak range $r\leq s$
    printed in Lemma~\texttt{B\_expand} of
    \cite{Molnar2018SemiInjective} is false: it introduces empty central
    words and omits the pure residual term.  Formula
    (\ref{eq:mps_reduction_residual_expansion}) is the corrected identity
    established in \cite{gap:mgsc18_residual_expansion_empty_word_error}.
\end{theorem}

\begin{proof}
    \leanok
    \uses{thm:mps_reduction_residual_sandwich,
        thm:mps_reduction_strict_window_range}
    Proceed by induction on the word and split its first letter as
    $B^i=N^i+WA^iV$.  The induction hypothesis supplies the strict windows
    in the remaining word.  A positive residual word between two selected
    $WA^iV$ blocks vanishes by
    Theorem~\ref{thm:mps_reduction_residual_sandwich}.  The remaining terms
    are the pure residual word and the recursively defined windows.  Theorem
    \ref{thm:mps_reduction_strict_window_range} identifies the latter with
    the explicit sum in (\ref{eq:mps_reduction_residual_expansion}).
\end{proof}

\begin{theorem}[Residual trace vanishing and nilpotence]
    \label{thm:mps_reduction_residual_algebra_nilpotent}
    \lean{MPSTensor.IsReduction.trace_evalWord_eq_trace_evalWord_add_trace_reductionResidual,
        MPSTensor.IsReduction.trace_evalWord_reductionResidual_eq_zero,
        MPSTensor.IsReduction.trace_eq_zero_of_mem_reductionResidualAlgebra,
        MPSTensor.IsReduction.isNilpotent_of_mem_reductionResidualAlgebra,
        MPSTensor.IsReduction.reductionResidualNilpotencyLength_isBound}
    \leanok
    \uses{def:same_mpv2_pos, def:mps_rectangular_reduction,
        def:mps_reduction_residual_nilpotency_length}
    Let $(V,W)$ be a reduction from $B$ to $A$, and assume
    $\mc{V}^+(B)=\mc{V}^+(A)$.  Then every nonempty residual word has trace
    zero, every element of $\mathcal R(B,A;V,W)$ has trace zero, and every
    element of this algebra is nilpotent.  The algebra is nilpotent, so the
    least residual-word nilpotency length $\ell(B,A;V,W)$ exists and is itself
    a bound.

    This is the nilpotency assertion in Proposition~21 of
    \cite{Molnar2018SemiInjective}.  The positive-length closed-chain
    equality is essential; a reduction alone does not imply nilpotency.  This
    is the $\lambda=1$ scope of the source argument; the dropped
    proportionality scalar is recorded in
    \cite{gap:mgsc18_reduction_proportionality_scalar}.
\end{theorem}

\begin{proof}
    \leanok
    \uses{thm:mps_reduction_residual_expansion,
        thm:mps_reduction_residual_exterior_sandwich,
        thm:mpu_shifted_zero_trace_power_nilpotent,
        thm:mpu_nil_matrices}
    Take the trace of the corrected residual expansion.  Cyclicity moves each
    mixed term to a positive residual word between $V$ and $W$, so it
    vanishes.  The unique term with no residual letters is $WA^{\mathbf i}V$;
    its trace is $\tr(A^{\mathbf i})$ because $VW=1$.  Thus
    \begin{align}
        \tr(B^{\mathbf i})
         & =\tr(A^{\mathbf i})+\tr(N^{\mathbf i}).
        \label{eq:mps_reduction_residual_trace_decomposition}
    \end{align}
    Equality of the two closed chains now gives
    $\tr(N^{\mathbf i})=0$, exactly as in
    \texttt{cornerproblem.tex} lines 3973--3981.  Linearity extends this
    identity to the generated nonunital algebra.  Vanishing traces of positive
    powers make each algebra element nilpotent, and finite-dimensional
    nil-algebra nilpotence supplies a residual-word bound.  Well-ordering gives
    the least such bound.
\end{proof}

\begin{theorem}[Dimension bound for the residual nilpotency length]
    \label{thm:mps_reduction_residual_dimension_bound}
    \lean{MPSTensor.IsReduction.reductionResidualAlgebra_list_prod_eq_zero,
        MPSTensor.IsReduction.bondDim_isReductionResidualNilpotencyBound,
        MPSTensor.IsReduction.reductionResidualNilpotencyLength_le_bondDim,
        MPSTensor.IsReduction.evalWord_reductionResidual_eq_zero_of_bondDim_le_length}
    \leanok
    \uses{thm:mps_reduction_residual_algebra_nilpotent,
        thm:mpu_nil_matrices}
    Under the hypotheses of the preceding theorem, every mixed product of
    $D_B$ elements of $\mathcal R(B,A;V,W)$ vanishes.  Consequently $D_B$ is
    a residual-word nilpotency bound,
    \begin{align}
        \ell(B,A;V,W) & \leq D_B,
        \label{eq:mps_reduction_residual_nilpotency_length_bound}
    \end{align}
    and $N^{\mathbf i}=0$ whenever $|\mathbf i|\geq D_B$.

    This numerical bond-dimension estimate is a derived TNLean/QICLean
    consequence of the finite-dimensional nil-matrix theorem.  It is not a
    statement of Proposition~21 in \cite{Molnar2018SemiInjective}.
\end{theorem}

\begin{proof}
    \leanok
    Apply the finite-dimensional nil-matrix theorem to the residual algebra.
    Its dimension-sharp matrix form makes every product of $D_B$ algebra
    elements zero.  Residual letters lie in this algebra, and prefix
    factorization then proves the three stated consequences.
\end{proof}

\begin{theorem}[Consequences of an arbitrary residual bound]
    \label{thm:mps_reduction_residual_bound_consequences}
    \lean{MPSTensor.IsReduction.evalWord_reductionResidual_eq_zero_of_bound_le_length}
    \leanok
    \uses{def:mps_reduction_residual_nilpotency_length}
    If $L$ is a residual-word nilpotency bound, then
    \begin{align}
        N^{\mathbf i} & =0
        \qquad\text{whenever } |\mathbf i|\geq L.
        \label{eq:mps_reduction_residual_bound_word}
    \end{align}
    This conclusion uses only the assumed bound; it does not require equality
    of closed chains.
\end{theorem}

\begin{proof}
    \leanok
    Split the residual word after its first $L$ letters.  The prefix is zero
    by the bound, so the full word is zero.
\end{proof}

\begin{theorem}[All-cut contracted residual identities]
    \label{thm:mps_reduction_all_cut_contracted}
    \lean{MPSTensor.reductionLeftContractedSum,
        MPSTensor.reductionRightContractedSum,
        MPSTensor.IsReduction.mul_evalWord_eq_reductionLeftContractedSum,
        MPSTensor.IsReduction.evalWord_mul_eq_reductionRightContractedSum}
    \leanok
    \uses{def:mps_rectangular_reduction,
        def:mps_reduction_residual_algebra}
    Let $(V,W)$ be a reduction from $B$ to $A$, and write
    $\mathbf i=(i_1,\ldots,i_n)$.  Then
    \begin{align}
        VB^{i_1}\cdots B^{i_n}
         & =\sum_{s=0}^{n}
        A^{i_1}\cdots A^{i_s}V
        N^{i_{s+1}}\cdots N^{i_n},
        \label{eq:mps_reduction_left_all_cut} \\
        B^{i_1}\cdots B^{i_n}W
         & =\sum_{r=0}^{n}
        N^{i_1}\cdots N^{i_r}W
        A^{i_{r+1}}\cdots A^{i_n}.
        \label{eq:mps_reduction_right_all_cut}
    \end{align}
    These identities require only the reduction.  They are precursors to the
    bound-truncated formulas in
    Theorem~\ref{thm:mps_reduction_contracted_expansions}.
\end{theorem}

\begin{proof}
    \leanok
    \uses{thm:mps_rectangular_reduction_identities,
        thm:mps_reduction_residual_sandwich,
        thm:mps_reduction_residual_expansion}
    Prove both identities by induction on the word.  For the left contraction,
    decompose the first letter and apply the corrected expansion to the tail;
    the sandwiched residual terms vanish and $VW=1$.  For the right
    contraction, decompose the first letter and apply the reduction identity
    to the tail.
\end{proof}

\begin{theorem}[Bound-truncated contracted residual expansions]
    \label{thm:mps_reduction_contracted_expansions}
    \lean{MPSTensor.reductionLeftContractedRangeSum,
        MPSTensor.reductionRightContractedRangeSum,
        MPSTensor.IsReduction.mul_evalWord_eq_reductionLeftContractedRangeSum,
        MPSTensor.IsReduction.evalWord_mul_eq_reductionRightContractedRangeSum}
    \leanok
    \uses{def:mps_rectangular_reduction,
        def:mps_reduction_residual_nilpotency_length}
    Let $(V,W)$ be a reduction from $B$ to $A$, let $L$ be any residual-word
    nilpotency bound, and write
    $\mathbf i=(i_1,\ldots,i_n)$.  Then
    \begin{align}
        VB^{i_1}\cdots B^{i_n}
         & =\sum_{n-L\leq s\leq n}
        A^{i_1}\cdots A^{i_s}V
        N^{i_{s+1}}\cdots N^{i_n},
        \label{eq:mps_reduction_left_contracted_expansion} \\
        B^{i_1}\cdots B^{i_n}W
         & =\sum_{0\leq r\leq\min(L,n)}
        N^{i_1}\cdots N^{i_r}W
        A^{i_{r+1}}\cdots A^{i_n}.
        \label{eq:mps_reduction_right_contracted_expansion}
    \end{align}
    Here $n-L$ denotes truncated subtraction, equivalently
    $\max(0,n-L)$ in integer notation.  The endpoints with a residual word of
    length exactly $L$ are retained as zero terms.  These are the two
    contracted formulas in Lemma~\texttt{B\_expand} of
    \cite{Molnar2018SemiInjective}, printed respectively at
    \texttt{cornerproblem.tex} lines 3997 and 3998.
\end{theorem}

\begin{proof}
    \leanok
    \uses{thm:mps_reduction_all_cut_contracted,
        thm:mps_reduction_residual_bound_consequences}
    Start from (\ref{eq:mps_reduction_left_all_cut}) and
    (\ref{eq:mps_reduction_right_all_cut}).  In the left sum, every omitted
    term has residual suffix length at least $L$.  In the right sum, every
    omitted term has residual prefix length at least $L$.  Equation
    (\ref{eq:mps_reduction_residual_bound_word}) makes all these terms zero.
\end{proof}

\begin{theorem}[Exterior-buffer identity]
    \label{thm:mps_reduction_exterior_buffer}
    \lean{MPSTensor.IsReduction.evalWord_mul_reduced_exterior_eq_evalWord_append}
    \leanok
    \uses{def:mps_rectangular_reduction,
        def:mps_reduction_residual_nilpotency_length}
    Let $(V,W)$ be a reduction from $B$ to $A$, and let $L$ be any
    residual-word nilpotency bound.  If the central word
    $\mathbf c$ is nonempty and the exterior words $\mathbf p,\mathbf q$
    satisfy
    \begin{align}
        L & \leq |\mathbf p|+1,
          & L                   & \leq |\mathbf q|+1,
        \label{eq:mps_reduction_exterior_buffer_bounds}
    \end{align}
    then
    \begin{align}
        B^{\mathbf p}WA^{\mathbf c}VB^{\mathbf q}
         & =B^{\mathbf p\mathbf c\mathbf q}.
        \label{eq:mps_reduction_exterior_buffer}
    \end{align}
    No closed-chain equality is assumed: only the reduction and the supplied
    residual bound enter this identity.  This is a derived blocking lemma in
    TNLean; it is not stated in Lemma~\texttt{B\_expand} or in its later uses
    in \cite{Molnar2018SemiInjective}.
\end{theorem}

\begin{proof}
    \leanok
    \uses{thm:mps_rectangular_reduction_identities,
        thm:mps_reduction_residual_sandwich,
        thm:mps_reduction_residual_expansion,
        thm:mps_reduction_residual_bound_consequences}
    First use induction along the right exterior word to show that a positive
    residual block preceded by $V$ vanishes once its length together with the
    right buffer reaches $L$.  Expand only the left exterior word
    $B^{\mathbf p}$ by
    Theorem~\ref{thm:mps_reduction_residual_expansion}.  The left buffer bound
    removes the pure residual term, while the preceding induction and
    (\ref{eq:mps_reduction_residual_sandwich}) remove the strict-window terms.
    Hence
    $B^{\mathbf p}N^iB^{\mathbf q}=0$ under the two buffer bounds.

    Separately, induction on a word $\mathbf a$ gives
    $VB^{\mathbf a\mathbf q}=A^{\mathbf a}VB^{\mathbf q}$ from the right
    buffer bound.  Write the nonempty central word as
    $\mathbf c=i\mathbf a$ and decompose the corresponding letter as
    $B^i=N^i+WA^iV$.  The residual term is zero by the first step, and the
    second induction identifies the remaining term with
    (\ref{eq:mps_reduction_exterior_buffer}).
\end{proof}

\begin{lemma}[Gauge invariance of block injectivity]
    \label{lem:gauge_invariance_block_injective}
    \lean{MPSTensor.isNBlkInjective_of_gaugeEquiv}
    \leanok
    \uses{def:gauge_equiv, def:l_blk_injective}
    If $A$ and $B$ are gauge equivalent and $A$ is $L$-block injective, then
    $B$ is $L$-block injective.
\end{lemma}

\begin{proof}\leanok
    \uses{lem:eval_word_gauge}
    Conjugation by an invertible matrix maps the span of the length-$L$ words
    of $A$ onto the span of the length-$L$ words of $B$. Since conjugation is
    an automorphism of the full matrix algebra, fullness of the first span
    implies fullness of the second.
\end{proof}

\begin{theorem}[Gauge transport of normality]
    \label{thm:gauge_invariance_normal}
    \lean{MPSTensor.isNormal_of_gaugeEquiv}
    \leanok
    \uses{def:gauge_equiv, def:normal}
    If $A$ and $B$ are gauge equivalent and $A$ is normal, then $B$ is normal.
\end{theorem}

\begin{proof}\leanok
    \uses{lem:gauge_invariance_block_injective}
    Apply Lemma~\ref{lem:gauge_invariance_block_injective} at a blocking length
    for which $A$ is block-injective.
\end{proof}

\begin{remark}\leanok
    \uses{lem:primitive_implies_irreducible, thm:burnside_matrix, thm:normal_from_primitive_posdef, thm:algspan_to_normal, thm:blocking_primitivity}
    The algebraic characterization of normality in
    Definition~\ref{def:normal} --- eventual full matrix span,
    equivalently eventual block injectivity --- is equivalent to the
    spectral characterization used in
    \cite[Definition~IV.1]{Cirac2021Matrix}: after TP normalization, the
    transfer map $\E_A$ is a primitive channel, equivalently irreducible
    with trivial peripheral spectrum. This equivalence is proved in
    \cite[Proposition~3]{Sanz2010Quantum}. In what follows, the
    directions of the equivalence are supplied by
    Lemma~\ref{lem:primitive_implies_irreducible},
    Theorem~\ref{thm:normal_from_primitive_posdef} for primitive normalized
    tensors with positive-definite fixed point, the Burnside route through
    Theorem~\ref{thm:algspan_to_normal} for the aperiodic algebraic
    criterion, and Theorem~\ref{thm:blocking_primitivity}. In particular, every
    subsequent use of ``normal'' here refers to this single notion;
    Definition~\ref{def:is_normal_canonical_form} employs the spectral
    formulation, but it describes the same class of tensors.
\end{remark}

\section{Canonical form}

\begin{definition}[Block-injective canonical form]\label{def:canonical_form}
    \lean{MPSTensor.CanonicalForm}
    \leanok
    \uses{def:mps_tensor, def:injective}
    Here, a \emph{block-injective canonical form} for a
    tensor consists of:
    \begin{enumerate}
        \item a number of blocks $r \ge 1$,
        \item bond dimensions $D_1, \ldots, D_r$,
        \item injective tensors
              $\{A_k^i\}_{i=0}^{d-1}$ with $A_k^i \in \MN{D_k}$
              for each block $k \in \{1,\ldots,r\}$,
        \item scaling factors $\mu_1, \ldots, \mu_r \in \C$.
    \end{enumerate}
    The associated tensor is the block-diagonal tensor
    \begin{align}
        A^i
         & = \bigoplus_{k=1}^{r} \mu_k A_k^i.
        \label{eq:mps_canonical_form}
    \end{align}
    This block-injective formulation of the canonical form is used in
    \cite{PerezGarcia2007Matrix}; see also~\cite{Cirac2021Matrix} for the
    normal canonical form.
\end{definition}

\begin{remark}\label{rem:cf_injective_vs_normal}\leanok
    Here we use the block-injective canonical form:
    each block is injective in the sense of
    Definition~\ref{def:injective}.
    This is stronger than the NT-block canonical form of
    \cite{Cirac2021Matrix}, where the blocks are only normal in the
    spectral sense.
    Chapter~\ref{ch:canonical} treats the normal canonical form needed
    for the blocking-based multi-block theory, while
    the periodic irreducible-form extension is kept separate. This section
    isolates the stronger block-injective form used in the injective and
    block-injective arguments.
\end{remark}

\begin{definition}[Block projection]\label{def:block_projection}
    \lean{Matrix.blockProjection}
    \leanok
    For a finite direct sum $\bigoplus_k \C^{D_k}$, the block projection $P_j$
    is the matrix that is the identity on the $j$-th summand and zero on all
    other summands.
\end{definition}

\begin{definition}[Block-diagonal matrix]\label{def:block_diagonal_matrix}
    \lean{Matrix.IsBlockDiagonal'}
    \leanok
    A matrix on $\bigoplus_k \C^{D_k}$ is block diagonal if it is a direct sum
    of square matrices, one on each summand.
\end{definition}

\begin{lemma}[Off-block characterization]\label{lem:block_diagonal_iff_off_block_zero}
    \lean{Matrix.isBlockDiagonal'_iff_offBlock_zero}
    \leanok
    \uses{def:block_diagonal_matrix}
    A matrix on $\bigoplus_k \C^{D_k}$ is block diagonal if and only if every
    entry from one summand to a distinct summand is zero.
\end{lemma}

\begin{proof}\leanok
    The forward direction follows from the definition of a direct sum of
    matrices. Conversely, if all off-block entries vanish, the diagonal blocks
    reconstruct the original matrix.
\end{proof}

\begin{lemma}[Commutation with block projections]
    \label{lem:block_diagonal_of_commutes_block_projection}
    \lean{Matrix.isBlockDiagonal'_of_commutes_blockProjection}
    \leanok
    \uses{def:block_projection, def:block_diagonal_matrix, lem:block_diagonal_iff_off_block_zero}
    If a matrix on $\bigoplus_k \C^{D_k}$ commutes with every block projection
    $P_j$, then it is block diagonal.
\end{lemma}

\begin{proof}\leanok
    \uses{lem:block_diagonal_iff_off_block_zero}
    Compare the $(i,j)$ block of the identity $XP_j=P_jX$ for $i\ne j$.
    The right multiplication by $P_j$ keeps that block, while the left
    multiplication by $P_j$ kills it. Hence every off-block entry is zero,
    and Lemma~\ref{lem:block_diagonal_iff_off_block_zero} applies.
\end{proof}

\begin{definition}[Block-diagonal tensor from blocks]\label{def:block_diagonal_tensor}
    \lean{MPSTensor.toTensorFromBlocks}
    \leanok
    \uses{def:mps_tensor}
    Given $r$ blocks with bond dimensions $(D_k)_{k=1}^r$,
    per-block tensors $\{A_k^i\}$,
    and scaling factors $(\mu_k)_{k=1}^r$,
    the associated \emph{block-diagonal tensor} is the tensor of
    total bond dimension $D = \sum_{k=1}^r D_k$ defined by
    \begin{align}
        A^i
         & = \bigoplus_{k=1}^{r} \mu_k A_k^i.
        \notag
    \end{align}
\end{definition}

\begin{lemma}[Word evaluation of a block-diagonal tensor]
    \label{lem:eval_word_block_diagonal_tensor}
    \lean{MPSTensor.evalWord_toTensorFromBlocks_eq_reindex_blockDiagonal}
    \leanok
    \uses{def:block_diagonal_tensor, def:eval_word}
    For a word $w$, the word evaluation of a block-diagonal tensor remains
    block diagonal:
    \begin{align}
        A^w
         & = \bigoplus_k \mu_k^{|w|} A_k^w.
        \label{eq:mps_block_word_eval}
    \end{align}
\end{lemma}

\begin{proof}\leanok
    Induct on the word and multiply block-diagonal matrices block by block;
    every letter contributes one scalar factor $\mu_k$ on the $k$-th block.
\end{proof}

\begin{theorem}[MPV decomposition]\label{thm:mpv_decomposition}
    \lean{MPSTensor.mpv_toTensorFromBlocks_eq_sum}
    \leanok
    \uses{def:block_diagonal_tensor, def:mpv, lem:eval_word_block_diagonal_tensor}
    The MPV of a block-diagonal tensor decomposes as
    \begin{align}
        \ket{V^{(N)}(A)}
         & = \sum_{k=1}^{r} \mu_k^N \ket{V^{(N)}(A_k)}.
        \label{eq:mps_mpv_decomposition}
    \end{align}
    Equivalently, for each configuration $\sigma$ one has
    \begin{align}
        V^{(N)}(A)_\sigma
         & = \sum_{k=1}^{r} \mu_k^N V^{(N)}(A_k)_\sigma.
        \notag
    \end{align}
\end{theorem}

\begin{proof}\leanok
    \uses{def:block_diagonal_tensor, def:mpv, lem:eval_word_block_diagonal_tensor}
    By Lemma~\ref{lem:eval_word_block_diagonal_tensor}, the full word evaluation
    is block diagonal with blocks $\mu_k^N A_k^w$ for a word $w$ of length $N$.
    Taking the trace of a block-diagonal matrix gives the sum of the block
    traces: $\tr(A^w) = \sum_k \mu_k^N \tr(A_k^w)$.
\end{proof}

\section{Blocking}

\begin{definition}[Blocked tensor]\label{def:blocked_tensor}
    \lean{MPSTensor.blockTensor}
    \leanok
    \uses{def:mps_tensor}
    Given a tensor~$A$ and a blocking length $L$,
    the \emph{$L$-blocked tensor} is the tensor with physical index set
    $\{0,\ldots,d{-}1\}^L$ (hence physical dimension $d^L$) and the
    same bond dimension~$D$, whose matrices are indexed by words
    $(i_1, \ldots, i_L) \in \{0,\ldots,d{-}1\}^L$:
    \begin{align}
        (A^{[L]})^{(i_1, \ldots, i_L)}
         & = A^{i_1} A^{i_2} \cdots A^{i_L}.
        \label{eq:mps_blocked_tensor}
    \end{align}
    Blocking coarse-grains $L$ neighbouring sites into one tensor: the inner
    virtual bonds are contracted, the $L$ physical legs merge into a single
    composite index, and the two outer bonds remain as the bond of $A^{[L]}$.
    \begin{center}
        % Formula/source: Papers/1606.00608/MPDO-22-12-17-2.tex,
        % lines 227--231, defines C^i=B^{i_1}\cdots B^{i_L} and draws the
        % blocked word in I_Blocking.png.
        % Ink-to-index map: i_1,...,i_L are the free physical indices;
        % the two outer wires are the blocked tensor's virtual indices.
        % Contraction set: every horizontal bond inside the dashed enclosure
        % is summed; the enclosure and its label add no contraction.
        % Type verdict: a rank-(L+2) word tensor, regarded as a rank-three
        % tensor after combining the L physical indices.
        % Boundary signature: semantic virtual-west=1 | virtual-east=1 |
        % physical-up=L | physical-down=0; the finite schematic emits 3 up legs.
        \tenkzkernel
        \begin{tenkz}[cols=4, boundary=open, physical=up]
            \tn[skin=box, ports={90:physical:$i_1$}]{A} & \tn[skin=box]{A} &
            \tn[skin=dots]{} & \tn[skin=box, ports={90:physical:$i_L$}]{A}
            \tnmark[form=enclosure, label pos=w]{(1,1) .. (1,4)}{$A^{[L]}$}
        \end{tenkz}
    \end{center}
\end{definition}

\begin{lemma}[Blocked word evaluation]\label{lem:eval_word_block}
    \lean{MPSTensor.evalWord_blockTensor}
    \lean{MPSTensor.blockedConfigEquiv}
    \lean{MPSTensor.ofFn_blockedConfigEquiv}
    \leanok
    \uses{def:blocked_tensor, def:eval_word}
    If $w = (J_1, \ldots, J_m)$ is a word in the blocked alphabet,
    where each $J_k$ is an $L$-tuple in $\{0, \ldots, d{-}1\}^L$,
    then concatenating the block words gives a word $\widetilde{w}$
    of length $mL$ in the original alphabet and
    $(A^{[L]})^w = A^{\widetilde{w}}$.
\end{lemma}

\begin{proof}\leanok
    \uses{lem:eval_word_append}
    Induct on the blocked word $w$ and split off the first block word.
    Lemma~\ref{lem:eval_word_append} turns concatenation of blocks into
    multiplication of the corresponding evaluations.
\end{proof}

\begin{lemma}[Blocking preserves MPV coefficients]\label{lem:mpv_block}
    \lean{MPSTensor.mpv_blockTensor_eq_mpv_blockedFlatConfig}
    \leanok
    \uses{def:blocked_tensor, def:mpv}
    For each basis state $\sigma = (\sigma_1, \ldots, \sigma_N)$ of
    the blocked system, where each $\sigma_k$ is an $L$-tuple in
    $\{0, \ldots, d{-}1\}^L$, there exists a basis state
    $\widetilde{\sigma}
        = (\widetilde{\sigma}_1, \ldots, \widetilde{\sigma}_{NL})$
    of the original system such that
    $V^{(N)}(A^{[L]})_\sigma
        = V^{(NL)}(A)_{\widetilde{\sigma}}$.
\end{lemma}

\begin{proof}\leanok
    \uses{def:blocked_tensor, def:mpv, lem:eval_word_block}
    Concatenate the block words and apply the blocked word evaluation
    lemma (Lemma~\ref{lem:eval_word_block}).
\end{proof}

\begin{theorem}[Blocking preserves same MPV]\label{thm:blocking_same_mpv}
    \lean{MPSTensor.SameMPV.blockTensor}
    \leanok
    \uses{def:same_mpv, def:blocked_tensor}
    If $A$ and $B$ generate the same MPV family, then for any
    blocking length~$L$, the blocked tensors $A^{[L]}$ and $B^{[L]}$
    also generate the same MPV family.
\end{theorem}

\begin{proof}\leanok
    \uses{lem:mpv_block}
    By Lemma~\ref{lem:mpv_block},
    each MPV coefficient of the blocked tensor $A^{[L]}$ equals
    a coefficient of $A$ at the corresponding concatenated
    configuration.
    The same holds for $B^{[L]}$ (with the same concatenated
    configuration).
    Since $A$ and $B$ agree on all configurations, so do
    $A^{[L]}$ and $B^{[L]}$.
\end{proof}

\section{MPV overlap}

\begin{definition}[MPV as Hilbert-space vector]\label{def:mpv_state}
    \lean{MPSTensor.mpvState}
    \leanok
    \uses{def:mpv}
    We regard $\ket{V^{(N)}(A)}$ as an element of $\C^{d^N}$
    indexed by configurations
    $\sigma \in \{0,\ldots,d{-}1\}^N$,
    with $\sigma$-component $V^{(N)}(A)_\sigma$ as in
    Definition~\ref{def:mpv}.
    This allows us to use the standard Euclidean inner product
    on $\C^{d^N}$.
\end{definition}

\begin{definition}[MPV inner product]\label{def:mpv_inner}
    \lean{MPSTensor.mpvInner}
    \leanok
    \uses{def:mpv_state}
    The \emph{MPV inner product} between two tensors $A$ and $B$
    at system size~$N$ is
    \begin{align}
        \braket{V^{(N)}(A)}{V^{(N)}(B)}
         & :=
        \sum_{\sigma \in \{0,\ldots,d{-}1\}^N}
        \overline{V^{(N)}(A)_\sigma} V^{(N)}(B)_\sigma,
        \label{eq:mps_inner_product}
    \end{align}
    i.e.\ the Euclidean inner product on $\C^{d^N}$
    (conjugate-linear in the first argument).
\end{definition}

\begin{definition}[MPV overlap]\label{def:mpv_overlap}
    \lean{MPSTensor.mpvOverlap}
    \leanok
    \uses{def:mpv}
    The \emph{MPV overlap} between two tensors $A$ (of bond
    dimension $D_1$) and $B$ (of bond dimension $D_2$)
    at system size~$N$ is
    \begin{align}
        O_{AB}(N)
         & :=
        \sum_{\sigma \in \{0,\ldots,d{-}1\}^N}
        V^{(N)}(A)_\sigma \overline{V^{(N)}(B)_\sigma}.
        \label{eq:mps_overlap}
    \end{align}
    Diagrammatically, the overlap contracts the MPV ring of $A$ against the
    conjugate ring of $B$ along their shared physical indices, both virtual
    rings closed by the trace:
    \begin{center}
        % O_{AB}(N) = \sum_\sigma V^{(N)}(A)_\sigma
        % \overline{V^{(N)}(B)_\sigma}: the ket row is the A ring and the
        % conjugate bra row is the conjugate B ring. Each vertical pairing
        % sums one shared physical index; the periodic closure closes both
        % virtual rings, the upper ring's return above the ladder and the
        % lower ring's below it, so neither return crosses the physical
        % pairings. Hence every index is contracted and the network is a
        % scalar. The bracket makes the N-site contraction schematic.
        % Boundary signature: (virtual-west=0 | virtual-east=0 |
        % physical-up=0 | physical-down=0).
        \tenkzkernel
        \begin{tenkz}[rows={ket, bra}, boundary=periodic]
            \tn[skin=box]{A} & \tn[skin=dots]{} & \tn[skin=box]{A} \\
            \tn[skin=box]{\overline{B}} & \tn[skin=dots]{} &
            \tn[skin=box]{\overline{B}}
            \tnmark[form=bracket]{(2,1) .. (2,3)}{$N\text{ sites}$}
        \end{tenkz}
    \end{center}
    This pairing is linear in the first MPV coefficient and conjugate-linear
    in the second. By Lemma~\ref{lem:overlap_eq_inner}, it is the complex
    conjugate of the Hilbert-space inner product from
    Definition~\ref{def:mpv_inner}.
\end{definition}

\begin{lemma}[Overlap equals conjugate inner product]\label{lem:overlap_eq_inner}
    \lean{MPSTensor.mpvOverlap_eq_star_mpvInner}
    \leanok
    \uses{def:mpv_overlap, def:mpv_inner}
    $O_{AB}(N)
        = \overline{\braket{V^{(N)}(A)}{V^{(N)}(B)}}$.
\end{lemma}

\begin{proof}\leanok
    \uses{def:mpv_overlap, def:mpv_inner}
    The overlap sums $V_\sigma \overline{W_\sigma}$
    while the inner product sums $\overline{V_\sigma} W_\sigma$.
    Complex conjugation swaps these.
\end{proof}

\begin{lemma}[Overlap determined by positive-length MPV]%
    \label{lem:overlap_eq_of_pos_mpv}
    \lean{MPSTensor.mpvOverlap_eq_of_pos_mpv_eq}\leanok
    \uses{def:mpv_overlap}
    If $V^{(N)}(A)_\sigma = V^{(N)}(A')_\sigma$ and
    $V^{(N)}(B)_\sigma = V^{(N)}(B')_\sigma$ for all $N > 0$
    and all $\sigma$, then $O_{AB}(N) = O_{A'B'}(N)$ for all
    positive~$N$.
\end{lemma}
\begin{proof}\leanok
    \uses{def:mpv_overlap}
    Direct pointwise rewriting: if the MPV coefficients agree for
    all positive chain lengths, the overlap sums are identical.
\end{proof}
